% Section: Abstract
% Format: Verdecchia & Bogner (2025) structured abstract
%
% EM-DASH PROHIBITION: Do not use --- or \textemdash. Use colons or semicolons.

\begin{abstract}
Large Language Models (LLMs) increasingly orchestrate tool chains that invoke
external services. When such a chain carries an opaque cryptographic
credential, such as a JSON Web Token (JWT), the credential must remain
byte-for-byte unchanged: a one-character deviation can cause downstream
authentication to fail. We report a controlled experiment across
LLM-mediated Model Context Protocol (MCP) tool chains ($N\!=\!48{,}000$ \revjb{runs}
across 16 frontier LLMs), \revjb{a second condition using production-scale credentials}
($N\!=\!1{,}600$), and an industrial case
study. Across the controlled runs, the JWT Defect/Failure Rate (DFR) was
18.0\% (range 0.5--99.8\%; 12.5\% excluding the Mistral Large~3 outlier).
Failures increased with credential length, \revjb{and the corrupted bytes were
concentrated in the JWT payload segment}. In the \revjb{production-realistic condition},
an explicit relay instruction left
a residual DFR of 19.9\% ($+1.9$~pp, $p\!=\!0.16$), showing the mitigation
does not eliminate the failure. To address this integrity risk, we propose the
Credential Interception Vault, a reference architecture that separates
deterministic credential handling from probabilistic LLM orchestration, uses
credentials by reference, and is model-agnostic. By removing credentials from
the LLM inference path, the architecture eliminates direct byte relay by
construction; \revjb{we illustrate this approach in a production fintech deployment,
suggesting it largely removes downstream authentication failure}.
\end{abstract}
